\section{Motivation}

PK11000 has been shown to have restorative effects on the Y220C mutant p53  \textit{in vivo} and \textit{in vitro} \parencite{bauer2SulfonylpyrimidinesMildAlkylating2016}. It has also shown restorative effects \textit{in silico}, using molecular dynamics simulations to study the p53 DNA binding domain bound to DNA \parencite{hanInsightsRationalDesign2022} and to study the full-length p53 bound to DNA \parencite{stetsonRestorationY220CP532025}.

However, the \textit{TP53} gene encodes at least \num{12} different proteins (``isoforms'') generated by \num{4} alternative translation start sites and \num{3} alternative terminal splicing patterns. Despite the importance of p53 in regulating many cellular fate-related pathways and the fact that these isoforms have been known about for several decades, they still remain largely unexplored. When a mutation occurs in the \textit{TP53} gene, it will be propagated to any p53 isoform expressed in the cell. This means that if a cell has the Y220C mutant \textit{TP53}, any isoform of p53 it expresses will also have the mutation. However, PK11000 binds to Cys182, which is conserved across all isoforms, so theoretically, it should at least be able to bind to each of the p53 isoforms. We may then ask a number of important questions, including but not limited to:

\begin{enumerate}
\item How do the isoforms differ from one another in terms of their dynamics?
\item How do p53-DNA interactions change as we make large N-terminal truncations or study alternative terminal splicing patterns?
\item Does the same pattern of disruption of the FL WT p53 dynamics extend to the other isoforms?
\item Does PK11000 even restore the mutant isoform dynamics? In what way does it or does it not? What can we learn about the mechanism of PK11000 rescue by looking at different isoforms?
\item Does the same pattern of PK11000 rescue of the FL Y220C p53 extend to the other isoforms?
\end{enumerate}

\section{Results}

\begin{figure}
\centering
\includegraphics[width=\textwidth]{Isoforms/Isoforms}
\caption[The p53 Splice Variants]{The twelve p53 isoforms. Each structure is from the final nanosecond of simulation, showing the Y220C\_PK11000 construct. Residues \numrange{1}{39} are shown in red, residues \numrange{40}{132} are shown in green, residues \numrange{133}{159} are shown in yellow, residues \numrange{160}{331} are shown in purple. The \(\alpha\) tail is shown in cyan. DNA is depicted in gray, a gray bead is the \ce{Zn}\textsuperscript{2+} ion, and PK11000 is shown in licorice.}
\label{fig:isoforms:all-structural}
\end{figure}

\begin{figure}
\centering
\includegraphics[width=\textwidth]{Isoforms/IsoformsOverview}
\caption[Structural Diagram of the p53 Isoforms]{A close-up of the WT \(\Delta 0\)p53\(\alpha\) isoform, or the full-length p53. Residues \numrange{1}{39} are shown in red, residues \numrange{40}{132} are shown in green, residues \numrange{133}{159} are shown in yellow, residues \numrange{160}{331} are shown in purple. The \(\alpha\) tail is shown in cyan. DNA is depicted in gray, and a gray bead is the \ce{Zn}\textsuperscript{2+} ion.}
\label{fig:isoforms:main-structural}
\end{figure}

To answer these questions, we have simulated each of the \num{12} splice variants bound to DNA in four forms: wild type p53 (WT), Y220C mutant p53 (Y220C), PK11000-bound wild type (WT\_PK11000), and PK11000-bound Y220C mutant p53 (Y220C\_PK11000).

When previously studying the FL protein, we considered the DNA binding domain (DBD) subset of the FL protein, defined from residues \numrange{96}{289} and denoted FL\textbar\textsubscript{DBD}, as defined by the 1TUP crystal structure of the p53 DNA binding domain \parencite{choCrystalStructureP531994}. This was useful because it allowed us to exclude from certain calculations the noisy N- and C-terminal domains, which made some patterns clearer. We would like an analogous subset of the p53 isoforms. However, isoforms with the \(\Delta133\) and \(\Delta160\) N-terminal deletions lack residues \numrange{96}{132} and \numrange{96}{160}, cutting into our previously defined subsets. We therefore define two additional subsets of the FL protein: the DBD', defined to encompass residues \numrange{133}{312}, and the DBD'', residues \numrange{160}{312}, denoted FL\textbar\textsubscript{DBD'} and FL\textbar\textsubscript{DBD''}. For isoforms with the \(\Delta160\) deletion, we say that there is no DBD' subset.

\figurepageref{fig:isoforms:all-structural} shows the twelve p53 isoforms in the Y220C\_PK11000 isoform, and \figurepageref{fig:isoforms:main-structural} shows a structural diagram of the \(\Delta0\)p53\(\alpha\) isoform.

\subsection{Protein RMSD Depends on Isoform and System}

\begin{figure}[htbp]
\centering
\includegraphics[width=\textwidth]{Isoforms/p53_Isoforms_FL_RMSD_vs_Time.png}
\caption[RMSD vs Time of p53 Isoforms]{Root Mean Square Deviation of protein structure over time. WT shown in blue, WT\_PK11000 shown in yellow, Y220C shown in red, and Y220C\_PK11000 shown in green. Averages over three replicates are shown in bold with a ghost representing the standard error of each frame.}
\label{fig:isoforms:rmsd-vs-time}
\end{figure}

We began by assessing our simulations for global structural stability, using RMSD. \figurepageref{fig:isoforms:rmsd-vs-time} shows RMSD as a function of time with each simulation's first frame as its reference.

RMSD increases from \num{0} at \qty{0}{\nano\second} (there is no difference between the starting structure and itself) and increases as the simulation proceeds. The majority of simulations have RMSD values that plateau around \qty{5}{\angstrom} and \qty{10}{\angstrom}.

RMSD values that steadily increase and are not bounded are typically indicative of a protein that has unfolded over the course of the simulation, whereas an increase in RMSD from a stable value usually indicates that the protein has adopted another conformation, one that is further from the starting structure by this distance metric. For instance, in isoform \num{3}, the Y220C\_PK11000 RMSD increases from approximately \qty{6}{\angstrom} to closer to \qty{10}{\angstrom}, but remains bounded at this value. This indicates that between \qtyrange{250}{500}{\nano\second}, the protein deviates from its conformation to a new conformation, but does not necessarily unfold. The Y220C system for the same isoform increases similarly shortly after the Y220C\_PK11000, indicating that it is undergoing a similar deviation from its starting structure. 

For the isoform \num{11} Y220C\_PK11000, we can see the RMSD for the first \qty{500}{\nano\second} is much greater than the last \qty{500}{\nano\second}. This is likely due to a local unfolding or an otherwise unusual conformational change of one replicate at the start of the simulation that returns to a more folded conformation towards the end. Figures \ref{fig:isoforms:rmsd-vs-time-DBDp} and \ref{fig:isoforms:rmsd-vs-time-DBDpp} on page \pageref{fig:isoforms:rmsd-vs-time-DBDp} show RMSD vs time for the FL\textbar\textsubscript{DBD'} and FL\textbar\textsubscript{DBD''} subsets of each isoform.

\begin{figure}
\centering
\includegraphics[width=\textwidth]{Isoforms/p53_Isoforms_RMSD_Violin.png}
\caption[p53 Isoform RMSD Violin Plot]{A violin plot showing, for each isoform and each system, the distribution of RMSD values. WT shown in blue, Y220C shown in red, and Y220C\_PK11000 shown in green. For each system, the median value of the distribution is shown with a horizontal line colored according to the system.}
\label{fig:isoforms:rmsd-violin}
\end{figure}

\figurepageref{fig:isoforms:rmsd-violin} shows the distribution of RMSD values as a violin plot for each isoform and system. Rather than showing RMSD as a function of time, a violin plot shows the distribution of the data. This allows us to identify different conformations that the system takes on by identifying RMSD values that are sampled more frequently during the simulation. Additionally, for each isoform, the median value of the distribution is shown with a horizontal line colored according to its system. This allows us to compare distributions of RMSD values between the mutated form of the protein with and without PK11000 bound, relative to the WT distribution.

In isoform \num{1}, we can see that the median value of the Y220C mutant is greater than that of the WT, but in the presence of the drug PK11000, the median value drops to below that of the WT. Isoform \num{8} is another such example. However, in isoforms \num{7} and \num{9}, we see very little difference between the median RMSD values of the mutant and drug-bound mutant proteins. Figures \ref{fig:isoforms:rmsd-violin-DBDp} and \ref{fig:isoforms:rmsd-violin-DBDpp} on page \pageref{fig:isoforms:rmsd-violin-DBDp} in Appendix \ref{chap:appendix:isoforms} show the RMSD violin plots for the FL\textbar\textsubscript{DBD'} and FL\textbar\textsubscript{DBD''} subsets of each isoform.

\begin{figure}
\centering
\includegraphics[height=0.4\textheight]{Isoforms/p53_Isoforms_RMSD_Diff.png}
\caption[p53 Isoform RMSD Difference Plot]{The \(\Delta\)RMSD values, computed between WT and Y220C and Y220C\_PK11000 and Y220C. For WT \textminus{} Y220C, a value less than zero means the Y220C protein deviates more from its starting structure than the WT; for Y220C\_PK11000 \textminus{} Y220C, a score less than zero means the presence of the drug lowered the average deviation from the starting structure.}
\label{fig:isoforms:rmsd-diff}
\end{figure}

To see more concisely how the presence of the Y220C mutation and the PK11000 restorative compound alters the relative protein stability, we computed the \(\Delta\)RMSD between the average RMSD values of the WT and Y220C p53 and between Y220C\_PK11000 and Y220C p53, shown in \figurepageref{fig:isoforms:rmsd-diff}. For \(\Delta\)RMSD\textsubscript{WT \textminus{} Y220C}, a value less than zero tells us that the mutant protein has a greater deviation from its starting structure than the WT protein. Similarly, for the \(\Delta\)RMSD\textsubscript{Y220C\_PK11000 \textminus{} Y220C}, a value less than zero means the presence of the drug lowers the average deviation from the starting structure.

For each isoform, we considered the \(\Delta\)RMSD of the full protein, and the FL\textbar\textsubscript{DBD'} and FL\textbar\textsubscript{DBD''} subsets. For the full protein, we can see that with the exception of p53\(\beta\), for isoforms \numrange{1}{8}, the presence of the Y220C mutation decreases the structural stability of p53 (seen via a greater deviation from the structure), and the addition of PK11000 is able to restore structural stability to mutant protein. For isoforms \numrange{10}{12}, the pattern is reversed: the Y220C mutant has a lower deviation from its starting structure compared to the WT, but the presence of PK11000 increases the flexibility of the protein back in the direction of the WT. 

Within our simulations, the disordered N- and C-terminal domains are attributable for the majority of the dynamics, as can be seen in the low RMSD values versus time in the FL\textbar\textsubscript{DBD'} and FL\textbar\textsubscript{DBD''} subsets of those proteins with large disordered domains (\(\Delta0\) and \(\Delta40\) N-terminal regions and the \(\alpha\) terminal splicing) in the supplementary data (Figures \ref{fig:isoforms:rmsd-vs-time-DBDp} and \ref{fig:isoforms:rmsd-vs-time-DBDpp} on page \pageref{fig:isoforms:rmsd-vs-time-DBDp} in Appendix \ref{chap:appendix:isoforms}). This is supported by small differences between systems for isoforms \numrange{1}{6}, as seen in the difference plot.

For the FL\textbar\textsubscript{DBD'} and FL\textbar\textsubscript{DBD''}, the \(\Delta\)RMSD values are too close to discern a large difference for most isoforms, except for isoforms \numrange{7}{9}, those with the \(\Delta133\) deletion. Isoforms \num{7} and \num{8} are destabilized by the Y220C mutant and re-stabilized by PK11000; isoform \num{9} is overly stabilized by Y220C and slightly destabilized by PK11000.

\begin{figure}
\centering
\includegraphics[height=0.4\textheight]{Isoforms/p53_RMSD_RescueScore.png}
\caption[p53 Isoform RMSD Rescue Score]{The PK11000 rescue score, using the metric of RMSD applied to the full protein, the DBD', and DBD''. A rescue score between 0 and 1 indicates a partial rescue, a score greater than 1 indicates an overshoot of the WT, and a score less than 0 indicates that Y220C\_PK11000 is further from the WT than Y220C.}
\label{fig:isoforms:rmsd-rescue-score}
\end{figure}

To see more quantitatively which simulations were rescued by PK11000 and by how much, we computed the RMSD rescue score for the three isoform subsets, shown in \figurepageref{fig:isoforms:rmsd-rescue-score}. When considering the entire protein, all isoforms, except for isoform \num{2} and \num{9}, show either partial rescue or an overshoot of rescue of the Y220C mutant. 

Given that isoforms \numrange{1}{9} have such similar RMSD values between the FL\textbar\textsubscript{DBD'} and FL\textbar\textsubscript{DBD''}, the rescue score for these subsets is nearly identical. Further, because isoforms \numrange{1}{6} have very similar RMSD values between WT, Y220C, and Y220C\_PK11000, the rescue score between the FL\textbar\textsubscript{DBD'} and FL\textbar\textsubscript{DBD''} is somewhat misleading. 

For the isoforms with the \(\Delta 133\) deletion, PK11000 can be seen to improve the structural stability of the FL\textbar\textsubscript{DBD'} of the protein, albeit slightly too much. Within the FL\textbar\textsubscript{DBD''} subset of isoforms \numrange{10}{12}, PK11000 is not seen to rescue isoforms \num{10} or \num{11}, though it does overshoot WT dynamics in isoform \num{12}.

\begin{figure}
\centering
\includegraphics[width=0.95\textwidth]{Isoforms/p53_Isoforms_Growing_RMSD.png}
\caption[Truncated p53 Isoform 1 RMSD Heatmaps]{Top: RMSD of truncated constructs of p53 isoform \num{1} as a function of truncation boundary. For each N-terminal starting residue \(n\), \(1 \leq n \leq 160\), and each C-terminal ending residue \(c\), \(312 \leq c \leq 393\), RMSD of the protein fragment between residues \(n\) and \(c\) was computed, then plotted as a heatmap. Bottom: For each N-terminal starting residue \(n\), the mean RMSD value averaging over all C-terminal ending residues \(c\).}
\label{fig:isoforms:rmsd-growing}
\end{figure}

\begin{figure}
\centering
\includegraphics[width=\textwidth]{Isoforms/Iso1_60-80}
\caption[Isoform 1 Residues that Increase RMSD]{The \(\Delta0\)p53\(\alpha\) with residues \numrange{60}{80} highlighted. These residues represent an extended region of the N-terminal domain with little secondary structure; including them in the calculation of RMSD increases it significantly.}
\label{fig:isoforms:rmsd-increasing-residues}
\end{figure}

Inspired by looking at how the RMSD of the protein changes with certain N-terminal deletions, we took the idea a step further. For each of the N-terminal residues \(n\) with \(1 \leq n \leq 160\), and for each of the C-terminal residues \(c\) with \(312 \leq c \leq 393\), we computed the average RMSD for the subset of isoform 1 starting at residue \(n\) and ending at residue \(c\), shown in \figurepageref{fig:isoforms:rmsd-growing}. The results were computed for each of the four systems, WT, Y220C, Y220C\_PK11000, and WT\_PK11000, and is shown as a heatmap with the x-axis corresponding to the N-terminal starting residue and the C-terminal ending residue of the computation.

It is immediately clear that there are certain residues in the N-terminal domain whose inclusion causes the RMSD to increase suddenly. To see this more clearly, we took the average RMSD value for each \(n\) over all the \(c\) values, \(312 \leq c \leq 393\). The inclusion of residues \numrange{60}{80} caused the RMSD to increase from \qty{4}{\angstrom} to \qty{8}{\angstrom} in all systems except for Y220C\_PK11000, which only increased a small amount. \figurepageref{fig:isoforms:rmsd-increasing-residues} shows the FL p53 with residues \numrange{60}{80} highlighted, and reveals that the region is a particularly flexible loop.

To compare this with the increasing RMSD with C-terminal inclusions, we repeated the same averaging, except for each \(c\), averaging over all \(n\) values, \(1 \leq n \leq 160\). We see that C-terminal inclusions have a lesser effect on the overall RMSD than N-terminal inclusions. Some tail-end fraying was observed, including the last \num{20} residues, which caused the RMSD to increase slightly, but this was to be expected. There were no subsets of residues in the C-terminal domain whose inclusion caused a large increase in the protein RMSD.

\subsection{Radius of Gyration Depends on Isoform and System}

\begin{figure}
\centering
\includegraphics[width=\textwidth]{Isoforms/p53_Isoforms_RoG_Violin.png}
\caption[p53 Isoform Violin Plot]{The radius of gyration for each isoform and each system with the distributions shown as a violin plot. WT is shown in blue, Y220C in red, and Y220C\_PK11000 in green.}
\label{fig:isoforms:rog-violin}
\end{figure}

The radius of gyration was computed to assess if the presence of the mutation or the restorative drug had an effect on the compaction of the protein. \figurepageref{fig:isoforms:rog-violin} shows the distribution of values as a violin plot. We can see that many of the isoforms have very similar radii of gyration across the wild type, mutant, and drug-bound mutant systems, namely isoforms \numlist{1;3;4;9}. In some cases, however, there is a difference. For example, in isoform \num{2}, the median radius of gyration for the mutant p53 is lower than that of the WT, indicating a compaction of structure upon mutation. When PK11000 is bound, we can see that the radius of gyration decreases further. An opposite pattern can be seen in isoform \num{5}, where the Y220C protein has a lower radius of gyration than the WT, and when PK11000 is bound to the mutant protein, we can see that the radius of gyration is restored to that of the wild type. Additionally, the violin plot reveals that in isoform \num{5}, the mutant protein has a small number of frames in which its structure has a compaction comparable to the median compaction of the WT, and similarly, the PK11000-bound Y220C protein has a small number of frames in which the protein has a conformation whose compaction is comparable with that of the WT, but a larger number of frames with a less compact conformation. This supports the idea that proteins must be considered by their ensemble of structures over time, rather than with a single canonical structure.

In a number of isoforms, namely isoforms \numlist{6;7;8}, we can see that there are a small number of frames in which the Y220C protein adopts a much less compact conformation with a greater radius of gyration. This pattern is not observed to such an extent in either the wild type or the drug-bound mutant, which provides another layer of computational evidence to support the observation that the mutant protein is less stable than the WT, and binding PK11000 restores some structural stability to the Y220C mutant protein.

In Appendix \ref{chap:appendix:isoforms}, the radius of gyration vs time is given for the FL, FL\textbar\textsubscript{DBD'}, and FL\textbar\textsubscript{DBD''} (Figures \ref{fig:isoforms:rog-vs-time}, \ref{fig:isoforms:rog-vs-time-DBDp}, and \ref{fig:isoforms:rog-vs-time-DBDpp}), in addition to violin plots for the FL\textbar\textsubscript{DBD'} and FL\textbar\textsubscript{DBD''} (Figures \ref{fig:isoforms:rog-vs-time-DBDp}, \ref{fig:isoforms:rog-vs-time-DBDpp}). Finally, the PK11000 rescue score is shown for the radius of gyration, \figurepageref{fig:isoforms:rog-rescue-score}. The rescue plot shows that, for the majority of isoforms, PK11000 either rescues or overshoots the WT-like radius of gyration values.

\subsection{RMSF Reveals Protein Core is Stable Across Isoforms}

\begin{figure}
\centering
\includegraphics[width=0.95\textwidth]{Isoforms/p53_Isoforms_RMSF.png}
\caption[p53 Isoform RMSF]{The root mean squared fluctuation of each of the isoforms by system over the course of the simulations. WT shown in blue, WT\_PK11000 shown in yellow, Y220C shown in red, and Y220C\_PK11000 shown in green. Horizontal gray lines are zero lines. Secondary structure is shown for each residue; blue indicates no secondary structure, red indicates a \(\beta\)-sheet for that residue, and yellow indicates an \(\alpha\)-helix.}
\label{fig:isoforms:rmsf}
\end{figure}

\begin{figure}
\centering
\includegraphics[width=0.95\textwidth]{Isoforms/p53_Isoforms_RMSF_Diff.png}
\caption[p53 Isoform RMSF Difference Plot]{The \(\Delta\)RMSF values, computed between WT and Y220C and Y220C\_PK11000 and Y220C. For WT \textminus{} Y220C, a value less than zero means the Y220C protein has more flexibility than the WT; for Y220C\_PK11000 \textminus{} Y220C, a score less than zero means the presence of the drug lowered the flexibility of that region. Secondary structure is shown for each residue; blue indicates no secondary structure, red indicates a \(\beta\)-sheet for that residue, and yellow indicates an \(\alpha\)-helix.}
\label{fig:isoforms:rmsf-diff}
\end{figure}

Whereas RMSD is a measure of global deviation from a reference structure and radius of gyration is a measure of protein compaction, RMSF is a measure of the flexibility of each residue in a protein over the whole simulation. \figurepageref{fig:isoforms:rmsf} shows the RMSF values for each residue of each isoform, separated by system. Even in globular proteins, the very ends of the protein may become slightly frayed, giving them a greater amount of flexibility than the core of the protein. In p53, which has long intrinsically disordered and N- and C-terminal domains, we see this effect amplified. The structured subset of p53 is typically defined between residues 96 and 289, and accordingly, within these regions, we can see that the RMSF values are quite low, indicating that there is low flexibility. Outside of this structured subset, there is little secondary structure, so the residues are much more flexible. In the \(\Delta133\) and \(\Delta160\) deletions, the entire disordered N-terminal region has been deleted, and there is no fraying. Within the structured subset of the protein, there are regions that have relatively higher flexibility: p53 is characterized as a \(\beta\)-sandwich, meaning it is comprised of multiple large \(\beta\)-sheets, and the regions that are observed to have higher flexibility correspond to bends or loops within this structure.

\figurepageref{fig:isoforms:rmsf-diff} plots the \(\Delta\)RMSF for each of the isoforms. To understand how the Y220C mutation and PK11000 alters local protein flexibility, we compute \(\Delta\)RMSF\textsubscript{WT \textminus{} Y220C} and \(\Delta\)RMSF\textsubscript{Y220C\_PK11000 \textminus{} Y220C}. For the \(\Delta\)RMSF\textsubscript{WT \textminus{} Y220C}, a value less than zero means that the Y220C protein is more flexible, and for the \(\Delta\)RMSF\textsubscript{Y220C\_PK11000 \textminus{} Y220C}, a value less than zero means the presence of the drug lowers the local flexibility in that region.

The difference plot reveals large differences between the WT, Y220C, and Y220C\_PK11000 structures in the N- and C-terminal domains. Broadly, in the N-terminal domain of isoforms \numrange{1}{3}, we see regions for which the Y220C mutation increases and decreases protein flexibility. In isoforms \numrange{4}{6}, those characterized by their \(\Delta40\) deletion, we can see that the Y220C mutation increases the flexibility. We can also see that the presence of PK11000 decreases the flexibility of residues \numrange{1}{95} in the Y220C mutant across most isoforms. There does not seem to be a clear pattern of behavior in the \(\alpha\) tail across isoforms. Given the little secondary structure character of both the N- and C-terminal domains, however, it is difficult to know if these findings are a true modulation of behavior or if the simulations simply behaved differently.

We can say more concretely, though, that in the structured subset of the protein, defined between residues \numrange{96}{289}, we see very little change between the three structures. 

\begin{figure}
\centering
\includegraphics[width=0.95\textwidth]{Isoforms/p53_Isoforms_RMSF_RescueScore.png}
\caption[p53 Isoform RMSF Rescue Score]{The rescue score of PK11000, computed using RMSF as the metric. A rescue score between 0 and 1 indicates a partial rescue, a score greater than 1 indicates an overshoot of the WT, and a score less than 0 indicates that Y220C\_PK11000 is further from the WT than Y220C. Secondary structure is shown for each residue; blue indicates no secondary structure, red indicates a \(\beta\)-sheet for that residue, and yellow indicates an \(\alpha\)-helix.}
\label{fig:isoforms:rmsf-rescue-score}
\end{figure}

To get a better understanding of whether PK11000 actually restored the WT dynamics of the protein, we computed the rescue score using RMSF as the metric. \figurepageref{fig:isoforms:rmsf-rescue-score} shows the rescue score of each residue for each of the isoforms. We observe that in isoforms \num{1} through \num{6}, many residues in the core domain are not assigned a rescue score, meaning that there was  little difference between their RMSF values. In the same isoforms, there are many regions in the N- and C-terminal domains in which the presence of the drug restores WT dynamics, overshoots WT dynamics, or worsens the mutant dynamics. Again, it is difficult to know whether these observations are real or if the disordered regions simply behave differently across the simulations. In isoforms \num{7} through \num{12}, with the exception of isoform \num{11}, we see that the rescue of the majority of residues has been overshot compared to that of the WT. This is likely related to our observations with RMSD and the radius of gyration, in which the presence of the drug also caused rescue to be overshot. The majority of residues in the \(\Delta160\)p53\(\beta\) structure are not rescued.

Comparing rescue with secondary structure, we can see that within the DNA-binding domain, regions that demonstrate any sort of change in dynamics by the Y220C mutant or PK11000 are often without secondary structure. This suggests that the those regions with secondary structure are more stable across constructs, and therefore do not see any disruption or restoration of the WT dynamics.

\subsection{Hydrogen Bonding Reveals PK11000-Dependent Restoration of WT Dynamics}

\begin{figure}
\centering
\includegraphics[width=\textwidth]{Isoforms/p53_Isoforms_HBonding_Metrics.png}
\caption[p53 Isoform Hydrogen Bonding]{Top: the Total Resonance Time (TRT, average number of hydrogen bonds between p53 and DNA) seen through several different lenses: all protein-DNA hydrogen bonds, all hydrogen bonds between the DNA and residues in the C-terminal domain (starting at residue 312), hydrogen bonds between DNA and residues K120, S241, R248, R273, A276, C277, R289, R283, and hydrogen bonds made between the DNA and K120, C277, and R280. Bottom: the Exhaustive Hydrogen Bond Sum (EHBS, total number of hydrogen bonds) seen through several lenses: all protein-DNA hydrogen bonds, all protein-DNA hydrogen bonds between the DNA and the C-terminal domain, all p53-DNA hydrogen bonds that persist for at least \qty{10}{\percent} of the simulation, and p53 CTD-DNA hydrogen bonds that persist for at least \qty{10}{\percent} of the simulation.}
\label{fig:isoforms:hbond}
\end{figure}

p53 is a transcription factor, and as such, one of its most important functions in the cell is to bind to the correct DNA sequences with the correct amount of strength to promote proper downstream gene transcription. For this reason, analyzing p53-DNA interactions, specifically how mutations and restorative compounds alter these interactions, gives us one of the most useful lenses to look at the behavior of p53 in the cell.

\figurepageref{fig:isoforms:hbond} begins with the total residence time. It is immediately apparent that those isoforms with the \(\alpha\) C-terminal tail have many more hydrogen bonds per frame than those with the \(\beta\) or \(\gamma\) C-terminal tails. This is confirmed with the second plot, which shows the total residence time of all of the hydrogen bonds made between the C-terminal domain of p53 and the DNA. We can see that these C-terminal interactions between the \(\alpha\) tail and the DNA are responsible for between roughly \num{3} and \num{5} additional hydrogen bonds per frame.

With regard to the specific and the nitrogenous-interacting residues, all except for lysine \num{120} are present across all isoforms; lysine 120 is missing from those isoforms with the \(\Delta133\) and the \(\Delta160\) deletions. There is no clear pattern between construct or isoform in the strength of the specific or nitrogenous-interacting hydrogen bonds, except for the marginal contribution that lysine \num{120} makes in those isoforms that contain it.

We then look to the EHBS to understand the specificity of the p53-DNA interactions. Every isoform across all systems makes between \num{400} and \num{900} unique hydrogen bonds over the course of the simulation. The \(\alpha\) tail contributes between \num{250} and \num{400} of these bonds. We can see by enforcing a cutoff of only \qty{10}{\percent} of frames, the total number of hydrogen bonds drops by a factor of \num{20} to a more reasonable \numrange{20}{40} hydrogen bonds, indicating the majority of these interactions are transient in nature and are therefore not representative of true p53-DNA interactions. Finally, we observe that the \(\alpha\) tail contributes roughly \numrange{6}{10} of these more significant interactions with the DNA.

\begin{figure}
\centering
\includegraphics[width=\textwidth]{Isoforms/p53_Isoforms_HBonding_Metrics_Diff.png}
\caption[p53 Isoform Hydrogen Bonding Difference Plot]{The \(\Delta\)TRT and \(\Delta\)EHBS for the lenses shown in \figurepageref{fig:isoforms:hbond}. For WT \textminus{} Y220C, a value less than zero means the Y220C protein makes more or stronger hydrogen bonds than the WT; for Y220C\_PK11000 \textminus{} Y220C, a score less than zero means the presence of the drug lowered the number or average strength of the p53-DNA hydrogen bonds.}
\label{fig:isoforms:hbond-diff}
\end{figure}

\begin{figure}
\centering
\includegraphics[width=\textwidth]{Isoforms/p53_Isoforms_HBonding_Metrics_RescueScore.png}
\caption[p53 Isoform Hydrogen Bonding Rescue Score]{The PK11000 rescue score when considering different hydrogen bonding metrics. For definitions of the metrics, see the caption of \figurepageref{fig:isoforms:hbond}. A rescue score between 0 and 1 indicates a partial rescue, a score greater than 1 indicates an overshoot of the WT, and a score less than 0 indicates that Y220C\_PK11000 is further from the WT than Y220C.}
\label{fig:isoforms:hbond-rescue-score}
\end{figure}

\figurepageref{fig:isoforms:hbond-diff} shows the \(\Delta\)TRT and \(\Delta\)EHBS plots for the lenses shown in \figurepageref{fig:isoforms:hbond}, and in \figurepageref{fig:isoforms:hbond-rescue-score}, we see the rescue score for each of these metrics.

In the full-length (FL) p53 (isoform \num{1}), we observed that the Y220C p53 makes weaker hydrogen bonds with the DNA compared to the WT, and that the presence of PK11000 increases the strength of these interactions beyond that of the WT \parencite{stetsonRestorationY220CP532025}. This can be seen with both \(\Delta\)TRT\textsubscript{WT \textminus{} Y220C} \(>0\) and \(\Delta\)TRT\textsubscript{Y220C\_PK11000 \textminus{} Y220C} \(>0\).  Accordingly, we see this reflected in our rescue score plot; PK11000 overshoots WT dynamics when bound to the full-length p53. Looking at the rescue plot, we can see that PK11000 modulates the TRT score in the direction of the WT in all isoforms except for isoform \num{4}, but comparing these results with the difference plot, we can see that the pattern of rescue differs between isoforms. Isoforms \numlist{5;10;11} follow the same pattern as the FL (Y220C makes hydrogen bonds weaker, PK11000 increases the strength), but in isoforms \numlist{2;6;8;9;12}, we see the opposite pattern: Y220C increases the average number of p53-DNA interactions, and PK11000 decreases the average number of hydrogen bonds. Both cases display a rescue of WT dynamics, but operating in opposite ways.

When considering the CTD TRT of isoforms with the \(\alpha\) tail (\(\Delta0\)p53\(\alpha\), \(\Delta40\)p53\(\alpha\), \(\Delta133\)p53\(\alpha\), \(\Delta160\)p53\(\alpha\)), we see that \(\Delta160\)p53\(\alpha\) follows the FL rescue pattern, isoform \(\Delta40\)p53\(\alpha\) follows the reverse FL rescue pattern, and \(\Delta133\)p53\(\alpha\) does not change upon Y220C mutation.

With the specific TRT, we can see in the FL that the Y220C mutation causes slightly stronger interactions between residues K120, S241, R248, R273, A276, C277, R289, R283, and the DNA. This same pattern can be seen in isoforms \numlist{8;9;12}, and the opposite pattern in isoforms \numlist{3;5;10;11}. For isoforms \numlist{2;4;6;7}, we do not see PK11000-induced rescue of the Y220C mutant.

Finally, for the nitrogenous-interacting residues, K120, C277, and R280, though there is no rescue in isoforms \num{1} through \num{3}, we do see a weakening of interactions relative to the WT in the Y220C mutant and a strengthening of interaction with PK11000 in isoforms \numlist{4;5;10;11}, and a strengthening of interactions in the Y220C mutant compared to the WT and a weakening of interactions back towards the WT by PK11000 in isoforms \numlist{6;7;8;9;12}.

We see similar patterns in the EHBS. In the FL p53 (isoform \num{1}), we observed the Y220C p53 making more total hydrogen bonds between p53 and the DNA, and PK11000 lowering this number in both the whole protein and the C-terminal domain. However, when we apply a cutoff, restricting our focus to hydrogen bonds that are present for only \qty{10}{\percent} of the simulation, we can see the Y220C mutation causes p53 to make fewer of these bonds, and that PK11000 restores the strength of the more significant interactions. We interpret this as the Y220C mutation causing p53 to make less specific interactions compared to the WT, and PK11000 restores the sequence specificity of the DNA interactions \parencite{stetsonRestorationY220CP532025}.

Extending this analysis to the different isoforms of p53 complicates the situation. Even disregarding the minimal CTD hydrogen bonds in isoforms with the short \(\beta\) or \(\gamma\) tails, isoform \num{11}, \(\Delta160\)p53\(\beta\), is the only other isoform that follows the same rescue pattern as the p53\(\alpha\). Isoform \num{8} sees a similar increase in total hydrogen bond count in the mutant form and a decrease when PK11000 is bound, but it also sees the same increase and decrease in the Y220C and Y220C\_PK11000 systems when the \qty{10}{\percent} cutoff is applied. Isoforms \numlist{9;10;12} see the opposite rescue pattern as the FL p53: their total hydrogen bond count is decreased in the presence of the Y220C mutation relative to the WT and increased upon PK11000 binding, and the number of hydrogen bonds present for at least \qty{10}{\percent} of the simulation increases in the Y220C mutant and then decreases in the Y220C\_PK11000 construct.

Interestingly, although isoform \num{11} lacked much structural restoration by the PK11000 drug, it showed rescue in every hydrogen bonding metric (excluding those analyzing CTD interactions).

\subsection{Local Energetic Interactions Are Restored by PK11000}

\begin{figure}
\centering
\includegraphics[width=\textwidth]{Isoforms/p53_Isoforms_EE.png}
\caption[p53 Isoform Embedding Error Score]{The local embedding error score for each residue for each isoform, relative to the WT. A greater value indicates a greater deviation from the WT. Red represents the difference between WT and Y220C, green represents the difference between WT and Y220C\_PK11000. Secondary structure is shown for each residue; blue indicates no secondary structure, red indicates a \(\beta\)-sheet for that residue, and yellow indicates an \(\alpha\)-helix.}
\label{fig:isoforms:ee}
\end{figure}

\begin{figure}
\centering
\includegraphics[width=\textwidth]{Isoforms/p53_Isoforms_EE_bar.png}
\caption[p53 Isoform Embedding Error Y220C vs Y220C\_PK11000]{A bar plot showing the average embedding error across each isoform. Red is the average EE score between WT and Y220C, green is the average EE score between WT and Y220C\_PK11000.}
\label{fig:isoforms:ee-bar}
\end{figure}

Our last analysis of the local energetic network around each residue. \figurepageref{fig:isoforms:ee} shows the embedding error score between the WT and both Y220C and Y220C\_PK11000 constructs for each residue of each isoform, a measure of similarity in local energetic structure relative to the WT. A greater value represents a greater deviation from the WT energetic interactions centered around that residue. Comparing across isoforms, we can see that the isoforms with fewer residues (isoforms \numrange{7}{12}) have a greater difference in energetic interactions compared to the WT. Within each isoform, there are a number of residues for which the Y220C is more similar to the WT than the Y220C\_PK11000, and where the Y220C\_PK11000 is more similar to the WT than the Y220C. \figurepageref{fig:isoforms:ee-diff} in Appendix \ref{chap:appendix:isoforms} shows the difference plot between Y220C\_PK11000 and Y220C EE scores.

To summarize the embedding error across each isoform, we took the average EE score for each isoform between the two constructs, and plotted it as a bar plot, shown in \figurepageref{fig:isoforms:ee-bar}. The difference is small for many isoforms, but for isoforms \numlist{1;4;7;8;9;10;12}, the Y220C\_PK11000 construct has a lower EE than the Y220C, indicating the protein's energetic interaction network is closer to the WT than the mutant. For \(\Delta133\)p53\(\beta\) and \(\Delta160\)p53\(\gamma\), the difference is particularly large, indicating a strong restoration in WT-like local energetic interactions across the protein.

\section{Discussion}

In this chapter, we have explored the restorative capabilities of the drug PK11000 across the twelve Y220C mutant p53 isoforms. We have done this by simulating each of the 12 isoforms across the wild-type, mutant, and drug-bound mutant constructs. We then used a variety of analyses to compare the dynamics across systems.

In terms of a restoration in global structural dynamics, when considering the entirety of each isoform, we observed either a partial restoration or overshoot of WT-like dynamics according to average RMSD from the starting structure in all isoforms except two and nine. The dynamics of the FL\textbar\textsubscript{DBD'} and FL\textbar\textsubscript{DBD''} subsets were very similar across all constructs studied. There was not a strong separation in radius of gyration values, even when considering the entirety of the protein, across most constructs.

When looking at local dynamics across the isoforms, as given by RMSF, we observed that the core domain of each isoform was very stable across the constructs. There was greater variation in the N- and C-terminal domains, though it was difficult to tell if these results were real or an artifact of their disordered nature. In the shorter isoforms, those with the \(\Delta133\) and \(\Delta160\) N-terminal truncations, we were able to observe a slight but widespread clamping of flexibility even in the core domain of the protein, and all isoforms except for isoform 11, \(\Delta160\)p53\(\beta\), showed either a partial rescue or an overshoot of WT-like dynamics.

An analysis of the hydrogen bonding between p53 and the DNA also showed a large amount of PK11000 rescue. In all isoforms except for isoforms four and eight, the presence of PK11000 overshot WT-like dynamics. There was little to no describable pattern in C-terminal hydrogen bonding with the DNA between isoforms with the \(\alpha\) tail, possibly suggesting that although the C-terminal domain is capable of making interactions with the DNA, they are not as important as those made by the DNA-binding domain.

Similar patterns of partial rescue or overshooting of WT-like dynamics can be seen in the specific TRT, nitrogenous-base interacting TRT, and EHBS measures, though interestingly, there is no relation between which isoforms are rescued for which metrics and which are not. Additionally, we see that isoform \num{11}, which lacked rescue in RMSF, showed rescue in RMSD and all relevant hydrogen bonding metrics. This suggests a deeper level of rescue than can simply be accounted for via a global metric or local metric.

Finally, studying the local energetic networks of the Y220C and Y220C\_PK11000 compared to the WT, the Y220C\_PK11000 showed a modest restoration of WT-like dynamics in most isoforms. This suggests that the binding of PK11000 is not simply affecting the local dynamics as a non-allosterically interacting small molecule might, but is actually modulating the electrostatic interaction network of all residues across the entire protein. 

In general, we hypothesize that the weight of the N- and C-terminal tails are able to modulate and limit the destabilizing effects of the Y220C mutation and therefore, little effect is seen with the binding of PK11000. This is evidenced not only by the larger isoforms having lower EE scores compared to smaller isoforms, but also in metrics such as RMSD of the FL\textbar\textsubscript{DBD'} subset of each isoform: the larger isoforms have significantly less variation compared to smaller isoforms.

In the end, we see that PK11000 was capable of modulating dynamics in the direction of WT-like dynamics for each mutant isoform of p53 (either via a partial rescue or overshoot of WT dynamics). This gives hope that, while it may not be clinically viable, PK11000 has general restorative capabilities and is not just limited to the FL p53 or the p53 DBD. It would have been surprising if it were limited to only these two p53 constructs, but nonetheless, proving this is an important step in reverse-engineering PK11000 to design new and better alternatives.

\section{Methods}

\subsection{Generation of Starting Structures}

Using the full-length (FL) p53 structure as a base \parencite{stetsonRestorationY220CP532025}, we generated the twelve different p53 splice variants. Isoform 1 is equivalent to the FL p53. The \(\Delta40\), \(\Delta133\), \(\Delta160\) N-terminal deletions were modeled by simply deleting the appropriate residues (residues \numrange{1}{39}, \numrange{1}{132}, or \numrange{1}{159}, respectively). The FL structure already encodes the canonical \(\alpha\) tail; C-terminal \(\beta\) and \(\gamma\) alternative splicings were generated by mutating residues \numrange{332}{341} and \numrange{332}{346} of the \(\alpha\) tail to the appropriate residues in the \(\beta\) and \(\gamma\) tails using the PyMol mutagenesis functionality \parencite{PyMOL} and trimming the excess residues.

\subsection{Simulation Protocol}
\label{chap:sim-protocol}

The isoform 1 trajectories were taken from \parencite{stetsonRestorationY220CP532025}. The WT and Y220C constructs of isoforms \num{2} through \num{12} were simulated by Fernando Caballero; for details regarding the simulation protocol, see \parencite{zotero-item-488}.

For the Y220C\_PK11000 and WT\_PK11000 simulations, three replicates were used, each \qty{1100}{\nano\second} long. To relax the engineered structures, we used an extended equilibration protocol. Eleven stages were used, with restraint weights equal to \numlist{100;75;65;55;45;30;20;15;10;5;1} \unit{\kilo\calorie\per\mol\per\angstrom^{2}}. The temperature was increased from \qty{0}{\kelvin} to \qty{300}{\kelvin} over \qty{20}{\pico\second} with \qty{20}{\kilo\calorie\per\mol\per\angstrom^{2}} restraint weights. Finally, a \qty{12}{\nano\second} equilibration phase was conducted.

All-atom simulations were performed with explicit solvent in an octahedral water box with periodic boundary conditions using the Amber 18 simulation suite \parencite{amber18}.  Particle-mesh Ewald electrostatics were used to handle long-range interactions. The ff14SB forcefield was chosen for protein dynamics \parencite{maierFf14SBImprovingAccuracy2015}, Bsc1 for DNA \parencite{ivaniParmbsc1RefinedForce2016}, and TIP3P for water \parencite{jorgensenComparisonSimplePotential1983}. The \ce{Zn}\textsuperscript{\num{2}+} atom was covalently bound to its coordinating residues, and parameterized using Zaff \parencite{petersStructuralSurveyZincContaining2010}. \ce{Na+} ions were added to maintain electroneutrality. Simulations were conducted using the GPU-optimized pmemd.cuda program \parencite{salomon-ferrerRoutineMicrosecondMolecular2013}. A timestep of \qty{1}{\femto\second} was used.

\subsection{Data Analysis}

Simulation data was first analyzed using CPPTRAJ \parencite{roePTRAJCPPTRAJSoftware2013}. The first \qty{100}{\nano\second} of simulation data was disregarded, and only the last \qty{1}{\micro\second} of data was analyzed. Subsequent analyses were performed using Python, making extensive use of Pandas, Numpy, and Matplotlib.

Pairwise residue interactions were generated with SAWNERGY, a Python package built in-house by Yehor Mishchyriak \parencite{Mishchyriak_SAWNERGY_A_Python}. Pairwise atom-by-atom interactions were generated, sampling one frame every \qty{10}{\nano\second}, then collapsed to a residue-by-residue level matrix with an edge weight threshold of \num{0.85} applied. A local thresholding protocol was additionally applied. The graph heat kernel was generated from the normalized Laplacian matrix \parencite{BenMasters}. Candidate diffusion times were sampled on a logarithmic scale from the Laplacian spectrum, and an appropriate diffusion time was selected by examining the total explained variance retained in downstream principal component analysis (PCA) across tested values. PCA was applied to the heat kernel of each frame to generate a low-dimensional residue-level embedding. The resulting embeddings were then collected by residue across all frames, and the Wasserstein distance was applied to determine the distance between distributions across constructs.

\subsection{Structure Visualization}

Structures were explored and modified using PyMol, and visualizations were generated using VMD \parencite{HUMP96, STON1998}.